\section*{Chapter \ref{ch:b3:lab-techniques-3} --- Lab Techniques III: Research Practice}

\begin{solution}{exo:b3:lab-techniques-3:1}
$(0.50/1013)^3 = \num{1.2e-10}$. Leaks, gas adsorbed on the glass and outgassing of grease and septa
leave far more than that; dry, hot glassware and good joints matter more than
extra cycles.
\end{solution}

\begin{solution}{exo:b3:lab-techniques-3:2}
$\qty{20.0}{mmol}/\qty{1.48}{mol/L} = \qty{13.5}{mL}$.
\end{solution}

\begin{solution}{exo:b3:lab-techniques-3:3}
$120 + 10 \times 1.00782503 + 55.93493633 - 0.00054858 = \qty{186.01264}{u}$. Error: $10^6 \times (186.0129 - 186.01264)/186.01264 = +1.4$ ppm,
well within 5 ppm.
\end{solution}

\begin{solution}{exo:b3:lab-techniques-3:4}
Primary: the journal article, the patent, the thesis. Secondary: the review, the
handbook of constants. Tertiary: the textbook.
\end{solution}

\begin{solution}{exo:b3:lab-techniques-3:5}
$M = \qty{194.0}{g/mol}$: C $96.0/194.0 = 49.48$\,\%, H $10.0/194.0 = 5.15$\,\%, N $56.0/194.0 = 28.87$\,\%. Differences 0.17, 0.07 and 0.16
points: within 0.4, the analysis supports the formula.
\end{solution}

\begin{solution}{exo:b3:lab-techniques-3:6}
$\qty{170.0}{mg}/\qty{212.0}{g/mol} = \qty{0.802}{mmol}$; $c = \qty{0.802}{mmol}/\qty{0.54}{mL} = \qty{1.5}{mol/L}$ (two significant figures, set by the
volume read).
\end{solution}

\begin{solution}{exo:b3:lab-techniques-3:7}
Means 73.0 and 77.0\,\%; both standard deviations $\sqrt{10/4} = 1.58$; $s_p = 1.58$. $t = 4.0/(1.58\sqrt{2/5}) = 4.0$, larger than
2.31: the improvement is significant at the 95\,\% level.
\end{solution}

\begin{solution}{exo:b3:lab-techniques-3:8}
$F = (0.80/0.30)^2 = 7.1 > 5.05$: the first method is significantly less precise.
\end{solution}

\begin{solution}{exo:b3:lab-techniques-3:9}
Buy small bottles with inhibitor, date each on receipt and on opening, store
closed, dark and cool; test for peroxides before any distillation or evaporation
and at fixed intervals after opening; dispose of bottles past their date, or
positive, through the hazardous-waste route without concentrating them.
\end{solution}

\begin{solution}{exo:b3:lab-techniques-3:10}
Hazards: hydrogen released (fire, pressure), the reactivity of NaH with water,
and a known exothermic decomposition of DMF with sodium hydride that can run away
on heating; DMF is also toxic to reproduction. Controls: substitute the solvent
(THF or 2-methyltetrahydrofuran) or the base; if kept, run small first, at low
temperature, adding NaH in portions with the temperature followed, in a hood
behind a shield, venting through a bubbler, with a cooling bath ready; written
procedure, a second person present, protective equipment last.
\end{solution}

\begin{solution}{exo:b3:lab-techniques-3:11}
Relative contributions: 0.50, 0.50, $0.010/20.0 = 0.05$, $0.010/15.0 = 0.067$ and 0.10\,\%; combined
$\sqrt{0.25 + 0.25 + 0.0025 + 0.0044 + 0.010} = 0.72$\,\%. The two integrals dominate: better signal-to-noise and
careful baseline and phase correction matter more than the balance.
\end{solution}

\begin{solution}{exo:b3:lab-techniques-3:12}
$Q = (89.0 - 82.0)/(89.0 - 80.6) = 0.83 > 0.71$: the test flags it. Before rejecting it, check the notebook
for a cause (a misread burette, an air bubble, a different batch of titrant); if one is
found, reject it and say why; if not, keep it, report the test, or repeat the
measurement.
\end{solution}

\begin{solution}{pb:b3:lab-techniques-3:1}
\textbf{1.} $\ce{FeCl2 + 2NaC5H5 -> Fe(C5H5)2 + 2NaCl}$.

\textbf{2.} Sodium cyclopentadienide is destroyed by oxygen and water, and iron(II)
chloride is oxidised to iron(III) in moist air.

\textbf{3.} $\qty{5.00}{g}/\qty{126.8}{g/mol} = \qty{0.0394}{mol}$; $\times \qty{185.8}{g/mol} = \qty{7.33}{g}$; yield $5.86/7.33 = 80$\,\%.

\textbf{4.} $(0.10/1000)^3 = 10^{-12}$ (in theory).

\textbf{5.} By heating the dimer, which undergoes a retro-Diels--Alder reaction, and
distilling off the volatile monomer; at room temperature the monomer dimerises
again by a Diels--Alder reaction, so it is kept cold and used at once.

\textbf{6.} Each portion releases hydrogen and heat: adding it slowly keeps the gas
flow and the temperature under control.

\textbf{7.} \qty{186.0126}{u}; error $+1.4$ ppm.

\textbf{8.} C $120.0/185.8 = 64.58$\,\%, H $10.0/185.8 = 5.38$\,\%; found 64.41 and 5.47: differences $-0.17$ and $+0.09$, within 0.4.

\textbf{9.} The ten hydrogens are equivalent by symmetry, and the rings rotate fast.

\textbf{10.} A single signal, all ten carbons being equivalent.

\textbf{11.} A single-crystal X-ray structure.

\textbf{12.} No: as an 18-electron complex (\cref{ch:b3:organometallic-bonding}) it is stable in air, unlike its
precursors.

\textbf{13.} Its signals do not overlap ferrocene's, it is a pure, non-volatile,
non-hygroscopic solid that can be weighed accurately, and it is inert towards the
sample.

\textbf{14.} Full relaxation of every signal between scans (a delay of several $T_1$),
a uniform excitation, enough signal-to-noise, and careful phase and baseline
correction.

\textbf{15.} Ferrocene $\ce{C10H10Fe}$ \qty{185.8}{g/mol}; 1,3,5-trimethoxybenzene $\ce{C9H12O3}$ \qty{168.0}{g/mol}.

\textbf{16.} $P_x = 3.942 \times (3/10) \times (185.8/168.0) \times (15.0/20.0) \times 0.999 = 0.980$: 98.0\,\%.

\textbf{17.} Relative $\sqrt{0.50^2 + 0.50^2 + 0.05^2 + 0.067^2 + 0.10^2} = 0.72$\,\%; absolute $0.72\,\% \times 98.0 = 0.7$ percentage point.

\textbf{18.} Yes, but weakly: a 2\,\% impurity of similar composition changes C and H by less
than the analysis can detect; qNMR is the more informative purity measurement.

\textbf{19.} Sodium hydride (hydrogen with water, fire), hydrogen evolved, THF and
cyclopentadiene (highly flammable; THF a \omterm{def:b3:lab-techniques-3:peroxide}{peroxide former}), hot cracking of the
dimer, iron(II) chloride (irritant), and ferrocene (flammable solid, harmful).

\textbf{20.} Sodium hydride: weighed as an oil dispersion under inert gas, added in
portions to a stirred, cooled solution, the gas vented through a bubbler, no water
near. THF: dry and peroxide-free from a column, used in the hood under nitrogen.

\textbf{21.} Under inert gas and with cooling, by slow addition of isopropanol, then
ethanol, then water, waiting each time until gas evolution stops.

\textbf{22.} A reversible one-electron wave, $\ce{Fe(C5H5)2+}/\ce{Fe(C5H5)2}$, with peaks about 59 mV apart at \qty{25}{\celsius}
(\cref{ch:b3:electrode-kinetics}); the couple is fast, stable and nearly insensitive to the
solvent, which makes it a convenient internal reference.

\textbf{23.} The notebook: date, scheme, amounts and sources, the \omterm{def:b3:lab-techniques-3:assessment}{risk assessment}, what
was done and seen, yields and the file names of the \omterm{def:b3:lab-techniques-3:notebook}{raw data}. The \omterm{def:b3:lab-techniques-3:literature}{supporting
information}: the full procedure and all characterisation data with copies of the
spectra.

\textbf{24.} So that others, and the authors years later, can reopen and re-process them
whatever software they have: findable, accessible, interoperable and reusable.

\textbf{25.} The qNMR purity of the ferrocene sample is $98.0 \pm 0.7$\,\% (standard uncertainty).
\end{solution}
